\begin{ouverture}[calcul]{2}{Nombres rationnels et fractions}{Partie 1 · Nombres et calculs}
{Simplifier · Multiplier · Inverser · Diviser}
\objectif{Simplifier une fraction et exprimer l’opposé d’un nombre rationnel.}
\objectif{Multiplier et diviser des fractions en utilisant l’inverse.}
\objectif{Résoudre des problèmes de fraction d’une quantité ou d’un reste.}
\prerequis{Fractions égales, comparaison, sommes et différences de fractions (Cinquième).}
\prerequis{Règles des signes du chapitre 1.}
\end{ouverture}
\pagestyle{otMain}\essentiel\fontsize{11}{15}\selectfont
\phantomsection\label{cours-2-1}
\essSection{1}{Des quotients d’entiers}
\begin{definition}Une fraction $\frac ab$ est le quotient de deux entiers $a$ et $b$, avec $b\ne0$ : c’est le nombre qui, multiplié par $b$, donne $a$, donc $b\times\frac ab=a$. Un nombre rationnel est un nombre égal au quotient de deux entiers relatifs. Il peut être entier ($\frac{12}4=3$), décimal ($\frac78=0{,}875$) ou ni l’un ni l’autre ($\frac13$).\end{definition}
\begin{propriete}[Simplifier]On ne change pas un quotient en multipliant ou en divisant son numérateur et son dénominateur par un même nombre non nul : $\frac{42}{56}=\frac{3\times14}{4\times14}=\frac34$. Simplifier, c’est diviser les deux termes par un diviseur commun.\end{propriete}
\begin{propriete}[Opposé d’un rationnel]$\frac{-a}{b}=\frac{a}{-b}=-\frac ab$. Par exemple $\frac43+\frac{-4}3=\frac{4-4}3=0$ : $\frac{-4}3$ est donc l’opposé de $\frac43$. On écrit en général le dénominateur positif.\end{propriete}
\begin{exemple}[Rappel de Cinquième]Pour additionner, soustraire ou comparer, on écrit les fractions avec un même dénominateur positif : $\frac34-\frac56=\frac9{12}-\frac{10}{12}=-\frac1{12}$. Ainsi $\frac34<\frac56$.\end{exemple}
\clearpage
\phantomsection\label{cours-2-2}
\essSection{2}{Multiplier, inverser, diviser}
\begin{propriete}Pour $b\ne0$ et $d\ne0$ : $\frac ab\times\frac cd=\frac{a\times c}{b\times d}$. Prendre une fraction d’un nombre, d’une quantité ou d’une fraction revient à multiplier : les $\frac23$ de $\frac45$ valent $\frac23\times\frac45=\frac8{15}$.\end{propriete}
\label{fig-cours-2-2-cours-fraction-fraction}\figureGeo{cours-fraction-fraction}

\begin{definition}L’inverse d’un nombre $a$ non nul est le nombre qui, multiplié par $a$, donne 1. On le note $\frac1a$ : c’est le quotient de 1 par $a$. Si $a\ne0$ et $b\ne0$, l’inverse de $\frac ab$ est $\frac ba$. Zéro n’a pas d’inverse.\end{definition}
\begin{propriete}Diviser par un nombre non nul revient à multiplier par son inverse : $\frac ab:\frac cd=\frac ab\times\frac dc$, avec $b$, $c$, $d$ non nuls. Exemple : $\frac54:\frac23=\frac54\times\frac32=\frac{15}8$.\end{propriete}
\begin{attention}$\frac12+\frac13=\frac56$ et non $\frac25$ : on n’additionne pas les dénominateurs. Les priorités s’appliquent aussi aux fractions : dans $\frac56-\frac23\times\frac94$, le produit passe d’abord.\end{attention}

\clearpage\rubrique{Découverte}
\begin{blocLivre}{ex-2-1}\begin{exercice}[1]{Des fractions égales}
Simplifier $\frac{42}{56}$ et $\frac{36}{-48}$, puis compléter $\frac{-3}{4}=\frac{\,?\,}{20}$.
\end{exercice}\end{blocLivre}
\begin{blocLivre}{ex-2-2}\begin{exercice}[1]{Entier, décimal ou non}
Les nombres $\frac{12}4$, $\frac7{20}$, $\frac13$ et $\frac{-18}6$ sont rationnels. Lesquels sont entiers ? décimaux ?
\end{exercice}\end{blocLivre}
\begin{blocLivre}{ex-2-3}\begin{exercice}[1]{L’opposé d’un quotient}
Calculer $\frac57+\frac{-5}7$. Que peut-on en déduire ? Écrire ensuite $\frac3{-8}$ avec un dénominateur positif.
\end{exercice}\end{blocLivre}
\begin{blocLivre}{ex-2-4}\begin{exercice}[1]{Additionner}
Calculer $\frac34-\frac56$.
\end{exercice}\end{blocLivre}
\begin{blocLivre}{ex-2-5}\begin{exercice}[1]{Multiplier}
Calculer $\frac{-4}{7}\times\frac{21}{10}$.
\end{exercice}\end{blocLivre}

\rubrique{Entraînement}
\begin{blocLivre}{ex-2-6}\begin{exercice}[2]{Une fraction d’une fraction}
Un rectangle est partagé en 4 colonnes et 3 lignes de cases identiques. On colorie les $\frac34$ du rectangle, puis les $\frac23$ de la partie coloriée. Quelle fraction du rectangle représente cette seconde partie ? Écrire le calcul.
\label{fig-ex-2-6}\figureGeo{fraction-rectangle}
\end{exercice}\end{blocLivre}
\begin{blocLivre}{ex-2-7}\begin{exercice}[2]{Trouver un inverse}
Donner l’inverse de $-\frac53$, puis calculer $\frac23:\left(-\frac53\right)$.
\end{exercice}\end{blocLivre}
\begin{blocLivre}{ex-2-8}\begin{exercice}[2]{Opposé ou inverse}
Donner l’opposé et l’inverse de 4, de $-\frac16$ et de $0{,}25$.
\end{exercice}\end{blocLivre}
\begin{blocLivre}{ex-2-9}\begin{exercice}[2]{Un calcul enchaîné}
Calculer $\frac56-\frac23\times\frac94$.
\end{exercice}\end{blocLivre}
\begin{blocLivre}{ex-2-10}\begin{exercice}[2]{Diviser une somme}
Calculer $\left(\frac34-\frac12\right):\frac58$.
\end{exercice}\end{blocLivre}

\rubrique{Problèmes}
\begin{blocLivre}{ex-2-11}\begin{exercice}[3]{Qui a bu le plus ?}
Inès boit les $\frac25$ d’une bouteille de jus. Hugo boit ensuite les $\frac23$ de ce qui reste. Qui a bu le plus ? Quelle fraction de la bouteille reste-t-il ?
\end{exercice}\end{blocLivre}
\begin{blocLivre}{ex-2-12}\begin{exercice}[3]{Lire une portion du reste}
Un réservoir contient 60 L. On utilise $\frac25$ de son contenu, puis $\frac13$ de ce qui reste. Combien de litres reste-t-il ?
\end{exercice}\end{blocLivre}
\begin{blocLivre}{ex-2-13}\begin{exercice}[3]{Retrouver une quantité}
Les $\frac7{12}$ d’un budget valent 175 euros. Quel est le budget total et quelle somme n’a pas encore été utilisée ?
\end{exercice}\end{blocLivre}
\begin{blocLivre}{ex-2-14}\begin{exercice}[3]{Une affirmation à discuter}
Tom affirme : « Diviser un nombre positif par une fraction le rend toujours plus petit. » Donner un contre-exemple et préciser quand le quotient devient plus grand.
\end{exercice}\end{blocLivre}
\begin{blocLivre}{ex-2-15}\begin{exercice}[3]{Trois parts successives}
Une somme est partagée : Lina reçoit $\frac14$ du total, Sami $\frac23$ du reste et Zoé les 36 euros restants. Trouver le total et les deux premières parts.
\end{exercice}\end{blocLivre}

\clearpage
